% reducesymm/QFT/stochastic.tex
% Predrag  excised from finiteQED.tex              		jul 27 2026
%%%%%%%%%%%%%%%%%%%%%%%%%%%%%%%%%%%%%%%%%%%%%%%%%%%%%%%%%%%%%%%%%%%%%%%%%%%

\chapter{Learning lattice quenched QED}
\label{c-stochastic}
%was \label{sect:latticeQED}

In 1924 Ryuichiro Kitano\rf{Kitano24} (for details, see \refpage{post:Kitano24})
applied in his
{\em QED Five-Loop on the Lattice},
% Progress of Theoretical and Experimental Physics, Volume 2025, 2025, 013B02, 
%Prog. Theor. Exp. Phys. 2025 013B02 (10 pages) DOI: 10.1093/ptep/ptae194
%\HREF{https://doi.org/10.1093/ptep/ptae194} {DOI}, 
\arXiv{2411.11554},
the lattice QED stochatic quantization method 
proposed in 
\refrefs{KiTaHa21,KitTak23}, \refpage{post:KiTaHa21} and  \refpage{post:KitTak23},
to evaluation of the quenched QED contribution of the {$g-2$} magnetic moment anomaly
to $\alpha^5$th (5-loop) order.

Kitano had performed his Monte-Carlo lattice simulations on five lattice 
\HREF{https://en.wikipedia.org/wiki/Hyperrectangle}{`hyperrectangular'} {\pcell}s,
\[
\BravCell{24^3}{48}{0}, 
\BravCell{28^3}{56}{0}, 
\BravCell{32^3}{64}{0}, 
\BravCell{48^3}{96}{0},
\BravCell{64^3}{128}{0}
\,, 
\]
in the notation of \toCLpdf{equation.74}{eq.~(76)},
and obtained values of the 4 and 5-loop contributions listed in table~\refeq{anomalValues}
that - though not competitive with the Feynman diagrams evaluations - were
in the right ball park.

Kitano, Takaura, and Hashimoto's  methodology is an example of
{\em numerical stochastic perturbation theory (NSPT)}, see the
notes of \refsect{sect:NSPT}.

%%%%%%%%%%%%%%%%%%%%%%%%%%%%%%%%%%%%%%%%%%%%%%%%%%
\section{Lattice quenched QED}
\label{sect:latticeQED}

We care about such nonperturbative QFT calculations, as he believes 
that at least QED is finite, in sense explained in \refref{Cvit77b}.

Who knew that the theory of deterministic chaos and turbulence 
turns out\rf{CL18} to be one of the corners of the standard stat mech and QFT,
the `strong coupling' regime of the Euclidean \spt\ field theory?
The spacetime translational invariance leads back to Euler, 
and replaces partion sums over \lsts\ (sums over `natural numbers') by zeta functions constructed from
\primeOrbs, products over `prime numbers' of a given deterministic theory.

The crux of the proposal is to replace Monte-Carlo simulations by
truncations of such zeta functions, each contribution given by a (numerically)
exact solutions of system's defining equations.

\subsection{A proposal: 
					  {\Dft} magnetic moment anomaly evaluation}
\label{sect:PCproposal}

\Spt\ theory of stochastic quantization has many moving parts.
A few initial steps are outlined in Domenico Lippolis paper \refref{Lippolis25}, \arXiv{2510.12532},
more in a paper under preparation.
\Dft\ is fundamentally different from Monte-Carlo evaluations of 
partition sums of stat mech and QFT, which it replaces by \spt\ zeta functions.
It is the yet unwritten 3rd millennium descendant of the quintessentially 20th 
century \wwwcb{} monograph.

\begin{enumerate}
	\item 
Determine  a {\steady} of QED deterministic {\ELe}.
If there is such, as in, for example,
\toCLpdf{section*.81}{appe.~C.2 \emph{Two-dimensional $\phi^4$ field theory \lst}}. In any space dimenion $d\geq1$.
	\item 
Determine  a small {\pcell} {\lst} of QED deterministic {\ELe}, for example a hyperrectangular \BravCell{2^3}{2}{0} \lst.
	\item 
Compute the {\stabexp}s of the above \lsts, as in 
\toCLpdf{section*.47}{sect.~X.B  \emph{{\Lst} stability}}.
	\item 
If you get that far, that's already pretty amazing. In that case,
automatize enumeration of {\Blatt}s 
\BravCell{\speriod{1}\times\speriod{2}\times\speriod{3}}{\period{}}{\tilt{}}, 
and determination of {\lsts} for each.
	\item  
	\begin{enumerate}
		\item  
Recompute  Kitano {$g-2$} magnetic moment anomaly,
now without Monte Carlo, but as a \detZeta\ sum over {\lsts}, using 
\toCLpdf{section*.59}{sect.~XI.F  \emph{{\Lsts} averaging formula}}.
		\item Can  you evaluate and bound the magnitue of gauge sets  of \refsect{sect:finitness}?
	\end{enumerate}
\end{enumerate}

%\bigskip
%%%%%%%%%%%%%%%%%%%%%%%%%%%%%%%%%%%%%%%%%%%%
\subsection{Lattice quenched QED notes}
\label{sect:latticeQEDnotes}

\begin{description}

\item[2025-04-08 Predrag].
\phantomsection \label{post:KiTaHa21}

Ryuichiro Kitano, Hiromasa Takaura, and 
\HREF{https://scholar.google.com/citations?hl=en&user=_P3ZbVgAAAAJ}
{Shoji Hashimoto}\rf{KiTaHa21}
{\em Stochastic computation of  {$g-2$}  in QED},
%J. High Energy Phys. 05, 199 (2021), 
\arXiv{2103.10106}:

``We perform a numerical computation of the anomalous magnetic moment {$g-2$} 
of the electron in QED by using the stochastic perturbation theory. 
Formulating QED on the lattice, we develop a method to calculate the 
coefficients of the perturbative series of {$g-2$} without the use of the 
Feynman diagrams. We demonstrate the feasibility of the method by 
performing a computation up to the $\alpha^3$  order and compare with the 
known results. This program provides us with a totally independent check 
of the results obtained by the Feynman diagrams, an example of the 
application of the numerical stochastic perturbation theory to physical 
quantities, for which the external states have to be taken on-shell. 
''

Reviews QED of {$g-2$}, but now on a Euclidean lattice.
They use using numerical stochastic perturbation theory (NSPT).

In the stochastic quantization, one obtains the expectation values
of operator products as fictitious time  $\tau$ averages of the field products. The
gauge field $A_\mu(n)$ is promoted to $A_\mu(n,\tau)$
governed by the Langevin equation
\begin{align}
 {\partial A_\mu(n,\tau) \over \partial \tau}
  = - {\delta S_{\rm lattice} \over \delta A_\mu(n,\tau)} + \eta_\mu (n,\tau).
 \label{KiTaHa21:eq:langevin}
\end{align}
The Gaussian noise $\eta_\mu (n,\tau)$ satisfies:
\[
 \langle \eta_\mu (n,\tau) \rangle_\eta = 0, \quad
 \langle \eta_\mu (n,\tau) \eta_\nu (n',\tau') \rangle_\eta
  = 2 \delta_{\mu \nu} \delta_{n n'}  \delta (\tau-\tau'),
\]
where $\langle \cdots \rangle_\eta$ denotes an ensemble average over
the random noise.
The ensemble average of field products converges to the correlation
functions obtained by the path integral quantization:
\[
 \langle A_{\mu_1} (n_1,\tau) \cdots A_{\mu_k} (n_k,\tau) \rangle_\eta \to
 \langle A_{\mu_1} (n_1) \cdots A_{\mu_k} (n_k) \rangle,
\]
as $\tau \to \infty$. This can be understood by the fact that the
probability distribution of the path integral, $e^{-S}$, is the fixed
point of the Fokker-Planck equation derived from the Langevin equation
in \refeq{KiTaHa21:eq:langevin}.
%
The ensemble average in the left-hand-side can be evaluated as the
``time'' average of the Langevin trajectories of the field product, i.e.,
\[
  \langle A_{\mu_1} (n_1) \cdots A_{\mu_k} (n_k) \rangle 
  = \lim_{\Delta \tau \to \infty} {1 \over \Delta \tau}
  \int_{\tau_0}^{\tau_0 + \Delta \tau} d\tau 
  A_{\mu_1} (n_1,\tau) \cdots A_{\mu_k} (n_k,\tau).
\]

Expand the field $A_\mu(n,\tau)$ as
\beq
 A_\mu(n,\tau) = \sum_{p=0}^{\infty} e^p A_\mu^{(p)} (n,\tau),
\ee{eq:expansion}
and solve the Langevin equation for each $A_\mu^{(p)}$.
%
The equation for gauge fields in each order $p$ is given by
\[
{\partial A_\mu^{(p)} (n,\tau) \over \partial \tau} 
= - {\delta S_{\rm lattice} \over \delta A_\mu(n,\tau)} \Bigg|_{(p)}
  + \eta_\mu (n,\tau) \delta_{p0},
\]
where $|_{(p)}$ denotes the $p$-th order term. The noise is applied
only for the lowest order, $A_\mu^{(0)}$. Higher order fields get
fluctuated through interactions with lower order fields.


For on-shell fermions, $s(p)^2 + m^2 = 0$ and $s(p+k)^2 + m^2 =
0$, the vertex function can be expressed by form factors as 
\bea
&& -i e_{\rm P} \bar u(p) \Gamma_\mu (p,k) u(p+k)
 =
\label{KiTaHa21:eq:form_factor}\\
&& \quad
-i e_{\rm P}  \bar u(p) \left(
F_1 (\hat k^2) \gamma_\mu
  - F_2 (\hat k^2)
 {\sigma_{\mu \nu} \hat k_\nu \over 2m}
 \right) u(p+k) +\mathcal{O}(a^2)
\,.
\nnu
\eea

Ward-Takahashi identities take familiar form
\beq
 \hat k_\mu G_\mu (p,k) 
= {e \xi \over \hat k^2} \left(
S (p+k) - S (p)
\right) e^{-2 \hat k^2 / \Lambda_{\rm UV}^2} .
\ee{KiTaHa21:eq:WT2}
Here $S(p)$ is the full fermion propagator, and $G_\mu (p,k)$ is
the vertex function.


\item[2025-04-08 Predrag].
\phantomsection \label{post:KitTak23}

Ryuichiro Kitano and Hiromasa Takaura\rf{KitTak23}
{\em Quantum electrodynamics on the lattice and
numerical perturbative computation of  {$g-2$} },
% Prog. Theor. Exp. Phys. 2023 103B02(24 pages)
%    journal = {Progress of Theoretical and Experimental Physics},
%               Prog. Theor. Exp. Phys. 2023, 103 (2023),
%    volume = {2023},
%   pages = {103B02},
%   (2023).
%     doi = {10.1093/ptep/ptad125},
%   url = {https://doi.org/10.1093/ptep/ptad125},
%   eprint =  {https://academic.oup.com/ptep/article-pdf/2023/10/103B02/52750764/ptad125.pdf},
% DOI: 10.1093/ptep/ptad125
\arXiv{2210.05569}

``We compute the electron g factor to the  $O(\alpha^5)$ order on the 
lattice in quenched QED. We first study finite volume (FV) corrections in 
various infrared regularization methods to discuss which regularization 
is optimal for our purpose. We find that in QEDL the FV correction to the 
effective mass can have different parametric dependences depending on the 
size of Euclidean time t and match the ‘naive on-shell result’ only at 
the very large t region, $t \gg L$. We adopt finite photon mass 
regularization to suppress FV effects exponentially and also discuss our 
strategy for selecting simulation parameters and the order of 
extrapolations to efficiently obtain the g factor. We perform lattice 
simulation using small lattices to test the feasibility of our 
calculation strategy. This study can be regarded as an intermediate step 
toward giving the five-loop coefficient independently of preceding 
studies.''


\item[2025-04-08 Predrag].
\phantomsection \label{post:Kitano24}

Ryuichiro Kitano\rf{Kitano24}
{\em QED Five-Loop on the Lattice},
% Progress of Theoretical and Experimental Physics, Volume 2025, 2025, 013B02, 
%Prog. Theor. Exp. Phys. 2025 013B02 (10 pages) DOI: 10.1093/ptep/ptae194
%\HREF{https://doi.org/10.1093/ptep/ptae194} {DOI}, 
\arXiv{2411.11554}.

A method to evaluate {$g-2$} on the lattice was proposed and developed in 
\refrefs{KiTaHa21,KitTak23}. The no-lepton-loop part of the 
calculation corresponds to ignoring the fermion determinant in the path 
integral. In QED, ignoring the fermion determinant leads to the free 
theory of the gauge field and thus the path integral can be  
performed by generating Gaussian noises. Also, the absence of the lepton 
loop means that there is no renormalization of $\alpha$, and thus the 
perturbative expansion can be formulated in terms of the bare parameter 
in the Dirac operator.

``
[...] obtain the perturbative coefficients of
$g-2$ in QED by using lattice simulations.
%
The calculation is quite simple as we only have a single diagram to
evaluate.
%
By concentrating on the diagrams without a lepton loop, the
formulation is further simplified as one can generate photon
configurations by the free theory.''

The correlation functions are 
obtained as statistical averages of quantities evaluated under each 
configuration. 
The QED without the fermion action is a free theory. The Euclidean
lattice action in the Feynman gauge can be written as
\begin{align}
    S_{\rm QED} = {1 \over 2}
    \sum_{n, \mu}
    \hat A_\mu(n) (-\nabla^2 + m_\gamma^2) \hat A_\mu(n).
    \label{Kitano24:eq:action_mod1}
\end{align}
We take the unit $a=1$ lattice spacing. 
%
The photon field, $\hat A_\mu(n)$, is a ultraviolet (UV) regularized
one defined by acting a differential operator,
\[
     \hat A_\mu(n) \equiv H \left( -{\nabla^2 \over \Lambda_{\rm
UV}^2} \right) A_\mu(n),
\]
with a function $H(x)$ to satisfy $H(x) \to 0$ for $x \gg 1$ and $H(x)
\to 1$ for $x \ll 1$.
%
In this study, we take a sharp UV cut-off,
\[
    H(x) = \left \{
        \begin{array}{ll}
            0, & x \ge 1,\\
            1, & x < 1.
        \end{array}
        \right.
\]
In QED, this cut-off is gauge invariant.

For the calculations of the correlation functions, we use the naive
Dirac operator on the lattice:
\[
    (D)_{nm}^{\alpha \beta} = m \delta_{nm} \delta_{\alpha \beta}
   + {1 \over 2} \sum_\mu \left[
   (\gamma_\mu)_{\alpha \beta} e^{i e A_\mu(n)} \delta_{n+\mu,m}
   - (\gamma_\mu)_{\alpha \beta} e^{-i e A_\mu(n-\mu)} \delta_{n-\mu,m}
   \right].
\]
% There are unwanted doublers in the propagators obtained by inverting
% this Dirac operator. We will take care of those unwanted modes later.
The parameter $m$ is the fermion mass and $e$ is the coupling constant. 
We formally expand every quantity in terms of $e$, and evaluate the 
statistical averages of each coefficient. In this way, at any stage of 
the calculation, we do not need a specific value of $e$. The coupling 
constant $e$ is not a parameter in the simulation. This is different from 
lattice simulations of nonperturbative dynamics where the continuum limit 
is taken by tuning the coupling constant to a critical value. Since there 
is no lepton loop in our study, $e$ is not renormalized. This means 
that $e$ can be identified as the physical coupling constant. 

They measure fermion–fermion–current three-point functions to obtain the 
form factors, in the position space in the Euclidean temporal 
direction and in the momentum space in the spatial direction. 

Each photon
configuration, $A_\mu(n)$, is generated according to the free
field theory. We evaluate the quantity inside the bracket of
\beq
  G_\mu (t) = \Big \langle \sum_{\bf p'}
  D^{-1} (t_{\rm sink}, t; {\bf p}, {\bf p}') \gamma_\mu
  D^{-1} (t, t_{\rm src}; {\bf p' + k}, {\bf p+k}) \Big \rangle
  ,
\ee{Kitano24:eq:three-point}
and take the ensemble average, $\langle \cdots \rangle$.
%
The locations $t_{\rm src}$, $t_{\rm sink}$ and $t$ are those of two
fermions and the current operator, respectively. We fix $t_{\rm src}$
and $t_{\rm sink}$ at some particular locations, and view the
correlation function as a function of $t$ for later purpose.


We take $L\to\infty$ and $T\to\infty$. Since we introduce 
the photon mass $m_\gamma$ as the IR regulator, the finite volume effects 
are suppressed exponentially as $e^{-m_\gamma L}$ and $e^{-m_\gamma T}$. 
Therefore, we keep the combinations of $m_\gamma L$ and $m_\gamma T$ 
large. 

 Projections to the electric and the magnetic functions are
\[
    G_E (t) = {\rm tr} \left[
        {1 + \gamma_4 \over 2} G_4 (t)
    \right], \quad
    G_M (t) = i \sum_{i,j,k} \epsilon_{ijk} {\rm tr} \left[
        {1 + \gamma_4 \over 2} 
        \gamma_5 \gamma_i G_j (t)
    \right] {\bf k}_k,
\]
where $i,j,k$ are indices for the spatial directions $x$, $y$ and
$z$. The Euclidean temporal direction is labelled as 4. The trace is
taken over the spinor indices. The Dirac representation of the gamma
matrices is convenient here since the projection, $(1+\gamma_4)/2$,
simply picks up the upper two components of the source and sink
spinors.

We perform the lattice simulations with five sets of lattice volumes,
$L^3 \times T$ = $24^3 \times 48$, $28^3 \times 56$, $32^3 \times 64$,
$48^3 \times 96$, and $64^3 \times 128$.

\end{description}

%%%%%%%%%%%%%%%%%%%%%%%%%%%%%%%%%%%%%%%%%%%%%%%%%%
\section{Numerical stochastic perturbation theory (NSPT)}
\label{sect:NSPT}
	
Kitano, Takaura, and Hashimoto's  methodology belongs to a large body of work which presumably has many
lessons relevant to, and implications for the proposed \dft\ calculation.

\begin{description}
\item[2026-07-26 Predrag] Have a look at:
Antonio González-Arroyo, Issaku Kanamori, Ken-Ichi Ishikawa, Kanata Miyahana, Masanori Okawa, Ryoichiro Ueno
\emph{Towards higher order numerical stochastic perturbation computation applied to the twisted Eguchi-Kawai model},
\arXiv{2001.02835}.

We have evaluated perturbation coefficients of Wilson loops up to $O(g^8)$ for the four-dimensional twisted Eguchi-Kawai model 
using the numerical stochastic perturbation theory (NSPT) in \arXiv{1902.09847}. 
We present the higher order calculation up to $O(g^{63})$.




\item[2026-07-26 Predrag] Have a look at:

Hironori Takei, Ken-Ichi Ishikawa, Masanori Okawa
\emph{The perturbative computation of the gradient flow coupling for the twisted Eguchi-Kawai model with the numerical stochastic perturbation theory},
\arXiv{2501.18175}.

The gradient flow method is a renormalization scheme in which the gauge field is flowed by the diffusion equation. 
The gradient flow scheme has benefits that the observables composed of flowed gauge fields do not require further renormalization and do not depend on the regularization. 
[...] We compute the gradient flow coupling for the twisted Eguchi--Kawai model using the numerical stochastic perturbation theory. 

\item[2013-10-17 Predrag, Kamal Sharma] 
Stochastic quantization with stochastic gauge fixing
introduces a gauge-fixing force which is tangent to the gauge orbit. More
work needs to be done in this method.


\end{description}

\bigskip

%%%%%%%%%%%%%%%%%%%%%%%%%%%%%%%%%%%%%%%%%%%%%%%%%%%
\noindent
Due to the \spt\ translational invariance, all contributions to zeta functions are multi-periodic.
\Lsts\ are computed on $d$\dmn\ tori, while their {\stabexp}s are evaluated over Brillioun zones.
Some literature on evaluating such \lsts:

\begin{description}

\item[2026-07-26 Predrag] Have a look at:
Margarita Garcia Perez, Antonio Gonzalez-Arroyo, Masanori Okawa
\emph{Perturbative analysis of twisted volume reduced theories},
\arXiv{1311.3465}.

[...] the perturbative expansion of SU(N) Yang-Mills theories defined on a d-dimensional torus of linear size l with twisted boundary conditions
[...] Numerical results in 2+1 dimensions provide a test of how these ideas extend into the non-perturbative regime.


\item[2026-07-26 Predrag] Have a look at:

Yimbo Ji
{\em Research on Numerical Stochastic Process Perturbation Theory for Calculating Higher-Order Perturbations of Matrix Models}
\HREF{https://hiroshima.repo.nii.ac.jp/records/2001821\#}{PhD thesis} (2024).
%\rf{Ji24thesis}

 I study the high-order perturbative behavior of a QCD-like 
 matrix  twisted reduced principal chiral model using \emph{numerical stochastic perturbation theory} (NSPT). 
 [...] This task proves challenging using conventional methods such as 
 Feynman diagrams, as the number of diagrams exhibits factorial divergence with a given perturbative order.
One way to alleviate this is NSPT which reduces the computational cost
from $O(N_{PT}!)$ to $O(N^2_{PT})$, where  $N_{PT}$ is the order in perturbation calculation
[...]. In this thesis,
I developed a new algorithm called the Paterson-Stockmeyer (P-S) method to accelerate the
bottleneck of the NSPT simulation. Moreover, using the newly developed method combined
with large N factorization, I have investigated the high-order perturbative behavior of TRPCM
and observed a signal indicating the existence of renormalons in the numerical results. In con-
clusion, I extracted non-perturbative information using high-order calculations of perturbative
coefficients for this model. The same workflow could be applied to other complicated
theories, such as full QCD.

Antonio González-Arroyo, Ken-Ichi Ishikawa, Yingbo Ji, Masanori Okawa
\emph{Perturbative study of large N principal chiral model with twisted reduction},
\arXiv{2208.01867}.
%doi 10.1142/S0217751X22502104

We compute the first four perturbative coefficients of the internal energy for the twisted reduced principal chiral model (TRPCM) using numerical stochastic perturbation theory (NSPT). [...]

\item[2026-07-26 Predrag] Have a look at:
Paolo Baglioni
\emph{Confronting Large Fluctuations in Numerical Stochastic Perturbation Theory}
\HREF{https://hdl.handle.net/1889/5700} {PhD thesis} (2024).
Tutore:
Prof. \HREF{https://www.pi.infn.it/~fdirenzo/} {Francesco Di Renzo}.
	
[...] Lattice gauge theories offer a framework for understanding non-perturbative aspects of quantum field theories. 
[...] INumerical Stochastic Perturbation Theory (NSPT)  
offers a fully automated stochastic method for calculating loop corrections in lattice field theories, 
to high perturbative orders. 
[...]
A not-well explored feature of NSPT is the freedom to choose any vacuum for calculating perturbation theory
[...] problems:  NSPT simulations exhibit large fluctuations in low-dimensional models. 
[...]
 In this thesis, we discuss NSPT simulations for the two-dimensional O(N) Non-Linear Sigma Models (NLSMs). 
[...]
we expect that huge fluctuations in simulations of low-dimensional models are somehow connected to the limited number of degrees of freedom. 
[...] Our numerical results show that in the large N limit NSPT simulations are not affected by the large fluctuations issue at high orders, 
[...]
 the larger the value of N, the more perturbative corrections we could compute
 [...]
 we discuss O(N) renormalon effects in the large N limit. 
 [...]
  we will provide new insights on the role of finite-volume effects.
  [...]
  Large N NSPT simulations for O(N) models can also be regarded as a preliminary step towards going back to perturbative expansions around non-trivial vacua. 
  
P. Baglioni, F. Di Renzo 
\emph{Large fluctuations in NSPT computations: a lesson from O(N) non-linear sigma models},
\arXiv{2402.01322}.

P. Baglioni, F. Di Renzo 
{Taming NSPT fluctuations in O(N) Non-Linear Sigma Model: simulations in the large N regime},
\arXiv{2412.02624} .

P. Baglioni, F. Di Renzo 
\emph{NSPT for O(N) non-linear sigma model: the larger N the better},
\arXiv{2401.11833}.

P. Baglioni, F. Di Renzo 
\emph{Numerical Stochastic Perturbation Theory around instantons},
\arXiv{2212.11240}.

Luigi Del Debbio, Francesco Di Renzo, Gianluca Filaci
\emph{Non perturbative physics from NSPT: renormalons, the gluon condensate and all that}
\arXiv{1811.05427}.

\item[2026-07-26 Predrag] Have a look at:
Georg Bergner, Antonio González-Arroyo, Ivan Soler
\emph{A $T_2×R^2$ roadmap to Confinement in SU(2) Yang-Mills theory},
\arXiv{2505.10396} .

We study the behaviour of $SU(2)$Y ang-Mills fields on a  $T_2×R^2$ geometry where the two-torus is equipped with twisted boundary conditions. We monitor the evolution of the dynamics of the system as a function of the torus size. For small sizes the behaviour of the system is well understood in terms of semiclassical predictions. In our case, the long distance structure is that of a two-dimensional gas of vortex-like fractional instantons with size and density growing with  torus size. Our lattice Monte Carlo simulations confirm the semiclassical predictions and allow the determination of the relevant scale signalling the transition to the non-dilute situation. 

%\item[2026-07-26 Predrag] Have a look at:
%
%\item[2026-07-26 Predrag] Have a look at:
%
%\item[2026-07-26 Predrag] Have a look at:


\item[2025-04-14 Predrag] Have a look at:

J. Dimock
{\em Quantum electrodynamics on the 3-torus},
\arXiv{math-ph/0210020}:
`` We study the ultraviolet problem for quantum electrodynamics on a 
three dimensional torus. We start with the lattice gauge theory on a 
toroidal lattice and seek to control the singularities as the lattice 
spacing is taken to zero. This is done by following the flow of a 
sequence of renormalization group transformations.'' 


\end{description}

%%%%%%%%%%%%%%%%%%%%%%%%%%%%%%%%%%%%%%%%%%%%%%%%%%%%%%%%%%%%%%%
\section{Beck notes}
\label{s:Beck}

\begin{description}
	\item[Predrag 2026-07-28]
Excerpts from 2025-08-31 siminos/spatiotemp/chapter/CML.tex
to be copied to here. Sometime "soon".
\end{description}




%%%%%%%%%%%%%%%%%%%%%%%%%%%%%%%%%%%%%%%%%%%%%%%%%%%%%%%%%%%%%%%%%%%%%%%%%
% \section{Method of smooth conjugacies}
% \label{sect:scfpo}
\input{scfpo}
%%%%%%%%%%%%%%%%%%%%%%%%%%%%%%%%%%%%%%%%%%%%%%%%%%%%%%%%%%%%%%%%%%%%%%%%%

%%%%%%%%%%%%%%%%%%%%%%%%%%%%%%%%%%%%%%%%%%%%%%%%
\printbibliography[heading=subbibintoc,title={References}]
